\documentclass{article}
\usepackage[UTF8]{ctex}
\usepackage{amsmath,amssymb,amsfonts}
\usepackage{enumitem}
\usepackage{geometry}
\geometry{a4paper, margin=1in}
% === 错题本排版优化（不参与题目标识） ===
\setlength{\parindent}{0pt}
\setlength{\parskip}{0.35em plus 0.10em minus 0.08em}
\setlength{\emergencystretch}{3em}
\setlist{itemsep=0.25em, topsep=0.40em, parsep=0pt}
\allowdisplaybreaks[2]
% === 错题本排版优化结束 ===
\begin{document}
\section*{中国矿业大学(北京)2025年数学分析真题}
\section*{26．中国矿业大学（北京）2025年数学分析试卷}

\subsection{题目1}
1．求极限 \(\lim _{n \rightarrow \infty} \frac{\int_{1}^{n} \ln \left(1+\frac{1}{\sqrt{x}}\right) \mathrm{d} x}{\sqrt{n}}=\)

\subsection{题目2}
2．曲面 \(x^{2}+y z+z^{2}-1=0\) 在点（1，1，0）处的切平面为 \(\_\_\_\_\)．

\subsection{题目3}
3．求不定积分 \(\int(x \ln x)^{\frac{3}{2}}(\ln x+1) \mathrm{d} x=\) \(\_\_\_\_\)．

\subsection{题目4}
4．求 \(\sum_{n=1}^{\infty} \frac{1}{n \cdot 2^{n}}=\) \(\_\_\_\_\)．

\subsection{题目5}
5．已知 \(L\) 为 \(x^{2}+y^{2}=2\) 的第一象限的部分取逆时针，求 \(\int_{L} x \mathrm{~d} y-2 y \mathrm{~d} x=\) \(\_\_\_\_\)．\\
\(\_\_\_\_\)

\subsection{题目6}
6．\(x=\frac{3}{4}\) 是 Riemann 函数的 \(\_\_\_\_\)点（填连续，可去间断，跳跃间断或第二类间断）．

\subsection{题目7}
7．计算 \(\lim _{x \rightarrow \infty}\left(\cos \frac{1}{x}+\sin \frac{2}{x}\right)^{x}=\) \(\_\_\_\_\)．

\subsection{题目8}
8．已知 \(\int_{0}^{+\infty} \frac{\sin x^{2}}{x^{p}} \mathrm{~d} x\) 绝对收敛，则 \(p\) 的取值范围为 \(\_\_\_\_\)．

二．（10 分）已知 \(D=\{(x, y) \mid x+y \leq 1, x \geq 0, y \geq 0\}\) ，且满足

\[
\lim _{x \rightarrow 1}\left(\frac{\ln x}{x-1}+\frac{a \sin x+\sin x}{x}\right)=\iint_{D} \cos \frac{x-y}{x+y} \mathrm{~d} x \mathrm{~d} y .
\]

求参数 \(a\) 的值．

三．（10 分）计算 \(\iint_{\sum} \frac{x \mathrm{~d} y \mathrm{~d} z+y \mathrm{~d} z \mathrm{~d} x+z \mathrm{~d} x \mathrm{~d} y}{\left(x^{2}+y^{2}+z^{2}\right)^{\frac{3}{2}}}\) ，其中 \(\sum\) 为 \(2 x^{2}+2 y^{2}+z^{2}=4\) 的外侧．

四．（10 分）假设 \(f(x)\) 在 \((-\infty,+\infty)\) 上一致连续，证明：存在非负实数 \(a, b\) ，使得对任意的 \(x \in (-\infty,+\infty)\) ，都有 \(|f(x)| \leq a|x|+b\) 。

五．（10分）叙述数列的 Cauchy 收敛准则，并用 Cauchy 收敛准则证明单调有界数列必收敛．

六．（10分）求函数 \(f(x)=x+\frac{1}{1+x}\) 的单调区间，凹凸区间，极值，拐点和渐进线．

\begin{itemize}
  \item[七．] （15 分）设 \(z=z(x, y)\) 由方程 \(9 x^{2}-54 x y+90 y^{2}-6 y z-3 z^{2}+18=0\) 所确定．
  \item[（1）] 求 \(\frac{\partial z}{\partial x}, \frac{\partial z}{\partial y}\) 。
  \item[（2）] 求 \(z(x, y)\) 的极值，极值点．
\end{itemize}

八．（15 分）设函数 \(f(x), g(x)\) 在 \([a, b]\) 上连续，在 \((a, b)\) 上可导，且 \(f(a)=f(b)=0, g^{\prime}(x)<0\) ．证明：至少存在一点 \(\xi \in(a, b)\) ，使得 \(\frac{f^{\prime}(\xi)}{g^{\prime}(\xi)}+\frac{\int_{a}^{\xi} f(t) \mathrm{d} t}{\int_{\xi}^{b} g(t) \mathrm{d} t}=0\) ．

九．（15 分）设数列 \(a_{n}, b_{n}\) 满足 \(e^{b_{n}}=e^{a_{n}}-a_{n}, n=1,2, \cdots\) ．

\begin{itemize}
  \item[（1）] 若 \(a_{n}>0\) ，证明：\(b_{n}>0\) ．
  \item[（2）] 若 \(a_{n}>0\) ，正项级数 \(\sum_{n=1}^{\infty} a_{n}\) 收敛，证明：\(\sum_{n=1}^{\infty} \frac{b_{n}}{a_{n}}\) 收敛。
\end{itemize}

\subsection{题目6}
原题附加题：（15 分）讨论下列函数列在给定区间是否一致收敛．

\begin{itemize}
  \item[（1）] \(f_{n}(x)=\arctan n x, x \in(0,+\infty)\) ．
  \item[（2）] \(f_{n}(x)=x \arctan n x, x \in(0,+\infty)\) ．
\end{itemize}
\end{document}
